\begin{center}
    \textbf{\Huge{Visão de Produto}}
\end{center}

\section{Introdução}

O propósito deste documento é coletar, analisar e definir as necessidades de alto nível e as principais \emph{features} do MIIA. O documento foca nas capacidades necessárias aos \emph{stakeholders} e usuários-alvo para viabilizar a interoperabilidade entre modelos de inteligência artificial (IA) treinados externamente e agentes autônomos executados em simuladores construtivos, com ênfase em cenários nos quais o simulador construtivo está integrado a ambientes de simulação virtual. Os detalhes de como o MIIA atende essas necessidades são especificados em artefatos complementares, incluindo casos de uso, \emph{user stories} e documentação suplementar.

\subsection{Referências}

\begin{itemize}
    \item Documento Institucional Preliminar, ``\textbf{Ambiente de Simulação Aeroespacial Avançado -- ASA$^{2}$}'', Instituto de Estudos Avançados (IEAv), 2025.
\end{itemize}

\newpage

\section{Posicionamento}

\subsection{Declaração do Problema}

\begin{table}[h]
    \begin{tabular}{|p{0.3\textwidth}|p{0.65\textwidth}|}
        \hline
            \textbf{O problema de} &
            Viabilizar a interoperabilidade, de forma controlada e determinística, entre modelos de inteligência artificial treinados externamente e agentes autônomos executados em simuladores construtivos, mantendo baixa latência, previsibilidade temporal e isolamento operacional, inclusive em ambientes nos quais o simulador construtivo está integrado à simulação virtual. \\
        \hline
            \textbf{afeta} &
            Engenheiros de simulação, integradores de sistemas, equipes de P\&D e operadores que dependem de agentes autônomos ou semi-autônomos capazes de atuar em ciclos de simulação. \\
        \hline
            \textbf{cujo impacto é} &
            A impossibilidade de empregar modelos de IA associados a agentes em simulações de nível operacional; acoplamentos frágeis, não determinísticos ou de difícil depuração; instabilidade do simulador; aumento de retrabalho e dívida técnica; e a incapacidade de validar cenários com rigor metodológico. \\
        \hline
            \textbf{uma solução bem-sucedida seria} &
            Um mecanismo padronizado, documentado e tecnicamente validado que viabilize a interoperabilidade entre diferentes modelos de IA e agentes autônomos, assegurando desempenho consistente, latência controlada, comportamento reproduzível e um processo rastreável e confiável de incorporação de modelos. \\
        \hline
    \end{tabular}
\end{table}

A solução envolve um conjunto de propriedades complementares:

\begin{itemize}

    \item \textbf{Determinismo:} para um mesmo conjunto de observações e estado interno, o modelo deve produzir exatamente as mesmas saídas, de forma consistente entre execuções. Isso elimina variabilidade indesejada e permite a reprodução rigorosa de cenários quando essas saídas são utilizadas pelos agentes.

    \item \textbf{Previsibilidade temporal:} o tempo necessário para que o modelo processe suas entradas e gere suas saídas deve permanecer dentro de limites bem definidos. A simulação depende de ciclos estáveis; qualquer variação temporal pode comprometer a execução do agente no ciclo do simulador.

    \item \textbf{Baixa latência:} a inferência do modelo deve ser suficientemente rápida para não atrasar o avanço do simulador. Em simulações de alta frequência, atrasos acumulados inviabilizam a operação dos agentes.

    \item \textbf{Isolamento operacional:} falhas, exceções ou comportamentos inesperados do modelo não podem comprometer o simulador nem os agentes em execução. A solução deve permitir detectar, conter e registrar falhas associadas aos modelos.

    \item \textbf{Reprodutibilidade:} as saídas produzidas por um modelo devem poder ser repetidas em diferentes execuções, máquinas ou versões do ambiente, desde que os mesmos artefatos (parâmetros e \emph{seeds}) sejam utilizados.

    \item \textbf{Rastreabilidade e verificação:} cada modelo deve possuir metadados claros (versão, origem, parâmetros e dependências). A execução deve permitir inspeção, instrumentação e validação contínua das inferências, facilitando a depuração e a avaliação metodológica.

\end{itemize}

Em conjunto, essas propriedades viabilizam a interoperabilidade controlada entre modelos de IA e agentes autônomos no ciclo de simulação.

\subsection{Declaração de Posicionamento do Produto}

\begin{table}[h]
    \begin{tabular}{|p{0.2\textwidth}|p{0.75\textwidth}|}
        \hline
            \textbf{Para} &
            Engenheiros de simulação, integradores de sistemas, equipes de P\&D e operadores que dependem de agentes autônomos ou semi-autônomos capazes de atuar de forma estável e previsível em ciclos de simulação. \\
        \hline
            \textbf{Que} &
            Enfrentam dificuldades técnicas ao empregar diferentes modelos de IA treinados externamente em agentes já implementados no simulador construtivo, resultando em acoplamento excessivo, dificuldades de validação e elevado custo de manutenção. \\
        \hline
            \textbf{O MIIA} &
            É um mecanismo especializado de interoperabilidade entre modelos de IA e agentes autônomos executados em simuladores construtivos. \\
        \hline
            \textbf{Que} &
            Viabiliza a interoperabilidade controlada entre modelos de IA e agentes autônomos, assegurando determinismo, baixa latência, reprodutibilidade e isolamento operacional no ciclo de simulação. \\
        \hline
            \textbf{Diferentemente de} &
            Servidores genéricos de inferência, arquiteturas distribuídas de propósito geral ou soluções \emph{ad-hoc} que introduzem forte acoplamento com a lógica do simulador e variabilidade indesejada no ciclo de execução. \\
        \hline
            \textbf{Nosso produto} &
           Prioriza desempenho temporal, estabilidade operacional, previsibilidade do comportamento dos agentes, rastreabilidade dos modelos e baixo acoplamento. \\
        \hline
    \end{tabular}
\end{table}

\subsection{Limites e Não-Objetivos do Produto}

O MIIA não se propõe a:

\begin{itemize}
    \item Integrar-se diretamente a simuladores virtuais ou substituir mecanismos de acoplamento entre simuladores construtivos e ambientes de simulação virtual, atuando exclusivamente no contexto do simulador construtivo.
    \item Garantir determinismo absoluto independentemente de variações de \emph{hardware}, compiladores, bibliotecas numéricas ou configurações de execução fora de um ambiente controlado.
    \item Substituir simuladores, motores físicos ou \emph{pipelines} de treinamento de agentes de IA.
    \item Corrigir deficiências intrínsecas dos modelos de IA, tais como instabilidade de políticas, vieses de treinamento ou comportamentos emergentes indesejados.
    \item Eliminar a necessidade de validação metodológica, testes experimentais e análise crítica dos resultados produzidos pelos modelos e agentes.
\end{itemize}

Esses limites são assumidos explicitamente como parte do escopo do produto.

\section{Descrição dos \emph{Stakeholders} e Usuários}

Esta seção descreve os principais \emph{stakeholders} e usuários envolvidos no projeto, bem como os problemas fundamentais percebidos no contexto atual da utilização de modelos de inteligência artificial por agentes autônomos em simuladores construtivos. O objetivo é fornecer o pano de fundo e a justificativa para a definição dos requisitos do sistema.

\newpage

\subsection{Resumo dos \emph{Stakeholders}}

\begin{table}[h]
    \begin{tabular}{| p{0.25\textwidth} | p{0.35\textwidth} | p{0.35\textwidth} |}
        \hline Nome & Descrição & Responsabilidades \\ \hline
        
        \textbf{Product Owner (PO)} &
        Responsável pela visão, prioridades e valor do produto &
        Define prioridades, valida entregas, aprova requisitos e delimita o escopo do produto \\ \hline
        
        \textbf{Pesquisador de Modelos de IA} &
        Responsável pelo desenvolvimento, treinamento e avaliação dos modelos de inteligência artificial &
        Fornece modelos treinados, documenta limitações, valida resultados e orienta decisões técnicas relacionadas às políticas de decisão \\ \hline

    \end{tabular}
\end{table}

\subsection{Resumo dos Usuários}

\begin{table}[h]
    \begin{tabular}{| p{0.30\textwidth} | p{0.30\textwidth} | p{0.30\textwidth} |}
        \hline Nome & Descrição & Responsabilidades \\ \hline
              
        \textbf{Engenheiro de Simulação} &
        Responsável pela configuração e manutenção do simulador construtivo &
        Utiliza o MIIA para estabelecer e manter a interoperabilidade entre modelos de IA e agentes autônomos, assegurando compatibilidade, estabilidade temporal e correto funcionamento no ciclo de simulação \\ \hline

    \end{tabular}
\end{table}

\subsection{Ambiente de Usuário}

Os usuários operam predominantemente em estações de trabalho Linux, integradas ao ambiente institucional do IEAv.

\subsubsection{Engenheiro de Simulação}

O Engenheiro de Simulação atua na configuração, operação e manutenção do simulador construtivo, sendo o principal usuário do MIIA. Suas atividades incluem a incorporação e substituição de modelos de inteligência artificial utilizados pelos agentes autônomos, a configuração de parâmetros de execução e a verificação do correto funcionamento do mecanismo no ciclo de simulação.

\subsection{Resumo das Principais Necessidades dos \emph{Stakeholders} e Usuários}

As necessidades a seguir refletem dois eixos complementares: (i) propriedades críticas de execução dos agentes no ciclo de simulação e (ii) requisitos de interoperabilidade, governança técnica e manutenção do produto ao longo do tempo.

\begin{table}[h]
    \begin{tabular}{| p{0.22\textwidth} | p{0.12\textwidth} | p{0.17\textwidth} | p{0.18\textwidth} | p{0.21\textwidth} |}
        \hline
        \textbf{Necessidade} & \textbf{Prioridade} & \textbf{Preocupações} & \textbf{Solução Atual} & \textbf{Soluções Propostas} \\ \hline
        
        Execução determinística dos agentes no ciclo de simulação &
        Alta &
        Variabilidade temporal, efeitos estocásticos dos modelos &
        Não há &
        Mecanismo padronizado para execução controlada de modelos com gestão explícita de estado e temporização \\ \hline
        
        Interoperabilidade simples e de baixo acoplamento entre modelos de IA e agentes &
        Alta &
        Dependências cruzadas, impacto no simulador &
        Não há &
        Interfaces estáveis e bem definidas entre modelos externos e agentes no simulador construtivo \\ \hline
        
        Reprodutibilidade entre execuções e ambientes &
        Alta &
        Sensibilidade a \emph{seeds} e dependências de execução &
        Não há &
        \emph{Pipeline} formal de exportação, validação e inspeção de modelos de IA \\ \hline
        
        Rastreabilidade e versionamento dos modelos &
        Média &
        Inconsistência de artefatos, perda de histórico &
        Não há &
        Processo padronizado de metadados, versionamento e documentação técnica dos modelos \\ \hline

        Previsibilidade temporal no ciclo de simulação &
        Média &
        Overhead de inferência e impacto perceptível na interação humana em cenários de simulação virtual &
        Não há &
        Execução sincronizada ao ciclo do simulador, compatível com simulações de alta frequência \\ \hline
        
    \end{tabular}
\end{table}

\subsection{Alternativas e Competição}

Diferentes abordagens técnicas podem ser consideradas para viabilizar a interoperabilidade entre modelos de inteligência artificial e agentes autônomos no ciclo de simulação construtiva, incluindo cenários nos quais o simulador construtivo encontra-se integrado a ambientes de simulação virtual. Essas abordagens representam alternativas tecnológicas viáveis, porém com diferentes compromissos em termos de desempenho temporal, determinismo, acoplamento e reprodutibilidade.

As principais classes de alternativas incluem:

\begin{itemize}
    \item \textbf{Runtimes interoperáveis de inferência} (ex.: ONNX Runtime): oferecem portabilidade entre modelos e ferramentas de treinamento, além de bom desempenho em inferência. Entretanto, carecem de mecanismos específicos para controle de estado, sincronização temporal e rastreabilidade no contexto de simulação construtiva, exigindo adaptações adicionais para garantir determinismo e previsibilidade operacional.

    \item \textbf{Runtimes nativos de inferência} (ex.: Libtorch, TensorFlow C++ API): permitem incorporação direta em aplicações C++ com menor \emph{overhead} e maior controle do ciclo de execução. Apesar do desempenho favorável, essas abordagens tendem a introduzir forte acoplamento com \emph{frameworks} específicos, além de maior complexidade de manutenção e validação ao longo do tempo.

    \item \textbf{Execução remota de inferência} (ex.: REST ou gRPC): proporciona flexibilidade arquitetural e isolamento de processos, facilitando a reutilização de modelos treinados externamente. Contudo, a comunicação remota introduz latência adicional, variabilidade temporal e dependência de infraestrutura de rede, tornando essa abordagem inadequada para ciclos de simulação de alta frequência que exigem previsibilidade rigorosa.

    \item \textbf{\emph{Embedding} de Python em aplicações C++}: possibilita acesso direto a modelos e ecossistemas de treinamento amplamente utilizados. No entanto, essa estratégia apresenta elevado acoplamento, dificuldades de depuração e riscos à estabilidade e à manutenção do simulador.
\end{itemize}

Embora cada uma dessas alternativas apresente vantagens em contextos específicos, nenhuma atende simultaneamente aos requisitos de determinismo, previsibilidade temporal, baixa latência, reprodutibilidade e baixo acoplamento exigidos pelo ambiente de simulação construtiva. O MIIA posiciona-se como uma solução especializada que busca endereçar essas limitações de forma coerente com o domínio de simulação construtiva.

\section{\emph{Overview} do Produto}

\subsection{Perspectiva do Produto}

O MIIA é um componente de infraestrutura especializado em prover interoperabilidade entre modelos de inteligência artificial e agentes autônomos executados em simuladores construtivos. Ele opera diretamente no ciclo de simulação, assegurando determinismo, previsibilidade temporal, rastreabilidade e uma interface padronizada entre modelos externos e agentes internos ao simulador.

O MIIA não substitui o simulador nem os \emph{pipelines} de treinamento de IA; atua como uma camada especializada responsável por garantir propriedades críticas de execução, alinhada às diretrizes institucionais para interoperabilidade padronizada de modelos de IA.

Esses mecanismos são empregados prioritariamente em ambientes nos quais o simulador construtivo encontra-se integrado a simuladores virtuais, impondo requisitos adicionais de previsibilidade temporal e estabilidade operacional.

\subsection{Suposições e Dependências}

\begin{itemize}
    \item Disponibilidade de modelos de IA treinados em PyTorch, SB3 ou ferramentas equivalentes.
    \item Execução predominante em ambiente Linux de alto desempenho.
    \item Requisitos operacionais de baixa latência e alta previsibilidade temporal no ciclo do simulador construtivo.
    \item Disponibilidade de especialistas técnicos em C++, simulação e inteligência artificial.
\end{itemize}

\section{\emph{Features} do Produto}

\begin{itemize}
    \item Interoperabilidade controlada entre modelos de IA e agentes autônomos no ciclo do simulador construtivo.
    \item Suporte a múltiplos formatos de exportação de modelos (ex.: PyTorch, ONNX).
    \item Mecanismos de rastreabilidade, versionamento e metadados dos modelos de IA.
    \item Instrumentação e monitoramento da inferência dos modelos durante a simulação.
    \item \emph{Pipeline} de validação mínima de modelos (integridade, compatibilidade e controle de \emph{seeds}).
\end{itemize}

\section{Outros Requisitos do Produto}

\begin{itemize}
    \item Latência compatível com ciclos de simulação construtiva de alta frequência, incluindo cenários de simulação virtual com interação humana.
    \item Robustez contra falhas no carregamento, inicialização ou execução de modelos de IA.
    \item API C++ simples, estável e de baixo acoplamento com o simulador construtivo.
    \item Documentação técnica cobrindo interoperabilidade, operação, validação e \emph{troubleshooting}.
    \item Garantias de segurança como requisito estrutural do produto, incluindo mecanismos para prevenir, detectar e conter a execução de modelos de IA maliciosos ou não conformes no contexto da simulação.
\end{itemize}
